\section{Complete Verification: Proof by Logical Evaluation} \label{sec:pbe}

Next, we formalize our
Proof By Logical Evaluation
algorithm \pbesym and show
that it is sound~(\S~\ref{sec:pbe:algo}),
that it is complete with respect to
equational proofs~(\S~\ref{sec:pbe:complete}),
and that it terminates~(\S~\ref{sec:pbe:terminates}).

\begin{figure}[t!]
\begin{tabular}{>{$}r<{$} >{$}r<{$} >{$}r<{$} >{$}l<{$} >{$}l<{$} }
\emphbf{Terms}    & \pred, \expr, \body  & ::=   & \smtlan\ \text{\kw{if}-free terms from Figure~\ref{fig:smtsyntax}}            & \\ [0.05in]

\emphbf{Functions}    & \Rfun  & ::=    & \rdef{x}{\pred}{\body}                   & \\ [0.05in]
\emphbf{Definitional Environment}  & \Renv  & ::=    & \emptyset \spmid \rfun \mapsto \Rfun, \Renv & \\ [0.05in]
\emphbf{Logical Environment} & \Penv  & ::=    & \emptyset \spmid \pred, \Penv               & \\ [0.05in]
\end{tabular}
\caption{{Syntax of Predicates, Terms and Reflected Functions.}}
\label{fig:pbe:syntax}
\end{figure}

\subsection{Algorithm}\label{sec:pbe:algo}

Figure~\ref{fig:pbe:syntax} describes
the input environments for \pbename.
%
The logical environment $\Penv$ contains
a set of hypotheses $p$, described in
Figure~\ref{fig:smtsyntax}.
%
The definitional environment $\Renv$
maps function symbols $f$ to their
definitions $\rdef{x}{\pred}{\body}$,
written as $\lambda$-abstractions
over guarded bodies.
%
Moreover, the body $\body$ and the
guard $\pred$ contain neither
$\lambda$ nor \texttt{if}.
%
These restrictions do not impact
expressiveness: $\lambda$s can be
named and reflected, and
\kw{if}-expressions can
be pulled out into top-level
guards using $\defIf\cdot$,
defined in~\citep{appendix}.
%
A definitional environment $\Renv$
can be constructed from $\RRenv$ as
%
$$
\tosmt{\RRenv} \defeq \{
   \gbind{f}{\lambda\overline{x}. \defIf{\tosmt{\refa}}} |
  (\gbind{f}{\lambda\overline{x}. \refa})\in\RRenv
  \}
$$


\begin{figure}[t!]
$$\begin{array}{lcl}
\toprule
\uinstsym & : & (\Renv, \Penv) \rightarrow \Penv \\
\midrule
\uinst{\Renv}{\Penv} & = &
  \Penv \cup \
  \bigcup_{\issubterm{\fapp{\rfun}{\expr}}{\Penv}}
    \inst{\Renv}{\Penv}{\rfun}{\exprs} \\[0.1in]

%\midrule
%\instsym & : & (\Renv, \Pred, \Rfun, \params{\Expr}) \rightarrow \setof{\Pred} \\
%\midrule
\inst{\Renv}{\Penv}{\rfun}{\exprs} & = &
  \left \{ \SUBSTS{\left(\tosmt{\fapp{\rfun}{x}} = \body_i\right)}{x}{\expr}\;
  \left|\;\, \left( \pred_i \Rightarrow \body_i \right) \in \overline{d}
  ,      \;\; \smtIsValid{\Penv}{\SUBSTS{\pred_i}{x}{\expr}}
  \right.
  \right \}\\
\quad \mbox{where} & & \\
\quad \quad \rbody{x}{d} & = & \Renv(\rfun) \\[0.1in]

\midrule
\pbesym & : & (\Renv, \Penv, \pred) \rightarrow \tbool \\
\midrule
\pbe{\Renv}{\Penv}{\pred} & = & \pbeloop{0}{\,\Penv\cup \
  \bigcup_{\issubterm{\fapp{\rfun}{\expr}}{\pred}}
    \inst{\Renv}{\Penv}{\rfun}{\exprs} } \\
\quad \mbox{where} & & \\
\quad \quad \pbeloop{i}{\Penv_i} & & \\
\quad \quad \spmid \smtIsValid{\Penv_i}{\pred}      & = & \ttrue \\
\quad \quad \spmid \Penv_{i+1} \subseteq \Penv_i    & = & \tfalse \\
\quad \quad \spmid \mbox{otherwise}                 & = & \pbeloop{i+1}{\Penv_{i+1}} \\
\quad \quad \quad  \mbox{where}                     &   & \\
\quad \quad \quad  \quad \Penv_{i+1}                & = & \Penv \cup \uinst{\Renv}{\Penv_i} \\
\bottomrule
\end{array}$$
\caption{\textbf{Algorithm \pbesym:} Proof by Logical Evaluation.}
\label{fig:pbe}
\label{fig:inst}
\label{fig:uinst}
\end{figure}

\mypara{Notation}
%
We write
$\issubterm{\rfun\ (\expr_1, \dots, \expr_n)}{\Penv}$
if the \smtlan term
$(\text{app} \dots (\text{app}\ f \ \expr_1) \dots \expr_n)$
is a syntactic subterm of some $\expr \in \Penv$.
%
We abuse notation to write
$\issubterm{\fapp{\rfun}{\expr}}{\expr'}$
for
$\issubterm{\fapp{\rfun}{\expr}}{\{\expr'\}}$.
%
We write \smtIsValid{\Penv}{\pred}
for SMT validity of the implication
$\Penv \Rightarrow \pred$.


\mypara{Instantiation \& Unfolding}
%
A term $\predb$ is a \emph{$(\Renv,\Penv)$-instance} if
there exists $\issubterm{\fapp{\rfun}{\expr}}{\Penv}$
such that:
%
\begin{itemize}
\item $\rlookup{\Renv}{\rfun} \equiv \rdef{x}{\pred_i}{\body_i}$,
\item $\smtIsValid{\Penv}{\SUBSTS{\pred_i}{x}{\expr}}$, and
\item $\predb \equiv \SUBSTS{(\fapp{\rfun}{x} = \body_i)}{x}{\expr}$.
\end{itemize}
%
A set of terms $Q$ is a \emph{$(\Renv, \Penv)$-instance}
if every $\predb \in Q$ is an $(\Renv, \Penv)$-instance.
%
The \emph{unfolding} of $\Renv, \Penv$ is the
(finite) set of all $(\Renv, \Penv)$-instances.
%
Procedure $\uinst{\Renv}{\Penv}$ shown in
Figure~\ref{fig:uinst} computes and returns
the conjunction of $\Penv$ and the unfolding
of $\Renv, \Penv$.
%
The following properties relate $(\Renv, \Penv)$-instances
to the semantics of \corelan and SMT validity.
Let \refapply{\RRenv}{e} denote the evaluation
of $e$ under the reflection environment $\RRenv$, \ie
$\refapply{\emptyset}{e} \doteq e$
and
$\refapply{(\RRenv, \bind{f}{e_f})}{e} \doteq \refapply{\RRenv}{\text{let rec } f = e_f \text{ in } e}$.

\begin{lemma}\label{lem:pbe:semantics}
For every $\env \models \RRenv$ and $\store\in \embed{\env}$,
\begin{itemize}
\item\textbf{Sat-Inst}  \label{lem:sat-inst}
If $\tosmt{e}$ is a $(\tosmt{\RRenv},\tosmt{\env})$-instance,
then \rsatisfies{\RRenv}{\store}{e}.
\item\textbf{SMT-Approx}\label{lem:smt-approx}
If $\smtIsValid{\tosmt{\env}}{\tosmt{e}}$,
then \rsatisfies{\RRenv}{\store}{e}.
\item\textbf{SMT-Inst} \label{lem:smt-inst}
If $\predb$ is a $(\tosmt{\RRenv},\tosmt{\env})$-instance
and $\smtIsValid{\extendpenv{\tosmt{\env}}{\predb}}{\tosmt{e}}$,
then $\rsatisfies{\RRenv}{\store}{e}$.
\end{itemize}
\end{lemma}

%%\begin{lemma}[\textbf{Sat-Inst}] \label{lem:sat-inst}
%%  If   $\satisfies{\Renv}{\store}{\Penv}$ and
%%       $Q$ is a $\Renv,\Penv$-instance,
%%  then $\satisfies{\Renv}{\store}{\Penv \wedge Q}$.
%%\end{lemma}
%%
%%\begin{lemma}[\textbf{SMT-Approx}] \label{lem:smt-approx}
%%If $\smtIsValid{\Penv}{\pred}$ then %forall $\store$,
%%   $\satisfies{\Renv}{\store}{\Penv}$ implies
%%   $\satisfies{\Renv}{\store}{\pred}$.
%%\end{lemma}
%%
%%\begin{lemma}[\textbf{SMT-Inst}] \label{lem:smt-inst}
%%If   (1)~$\satisfies{\Renv}{\store}{\Penv}$,
%%     (2)~$Q$ is an $\Renv,\Penv$-instance,
%%     (3)~$\smtIsValid{\Penv \wedge Q}{\pred}$
%%then $\satisfies{\Renv}{\store}{\pred}$.
%%\end{lemma}

\mypara{The Algorithm}
%
Figure~\ref{fig:pbe} shows our proof search algorithm
$\pbe{\Renv}{\Penv}{\pred}$ which takes as input
%
a set of \emph{reflected definitions} $\Renv$,
%
an \emph{hypothesis} $\Penv$, and
%
a \emph{goal} $\pred$.
%
The \pbename procedure recursively \emph{unfolds}
function application terms by invoking $\uinstsym$
until either the goal can be proved using the
unfolded instances (in which case the search
returns $\ttrue$) \emph{or} no new instances
are generated by the unfolding (in which case
the search returns $\tfalse$).

\mypara{Soundness}
%
First, we prove the soundness of \pbesym. %(and hence, our type system.)

\begin{theorem}[\textbf{Soundness}] \label{thm:pbe:sound}
If   $\pbe{\tosmt{\RRenv}}{\tosmt{\env}}{\tosmt{e}}$
then $\forall \store\in\embed{\env}$,
\evalsto{\apply{\store}{\refapply{\RRenv}{e}}}{\etrue}.
% we have $\satisfies{\Renv}{\store}{\Penv \Rightarrow \pred}$.
\end{theorem}

We prove Theorem~\ref{thm:pbe:sound}
using the Lemma~\ref{lem:pbe:semantics} % ,~\ref{lem:smt-approx},~\ref{lem:smt-inst}
that relates instantiation, SMT validity,
and the exact semantics.
%
Intuitively, \pbesym is sound as it
reasons about a finite set of instances
by \emph{conservatively} treating all
function applications as \emph{uninterpreted}~\cite{Nelson81}.

\subsection{Completeness} \label{sec:pbe:complete}

Next, we show that our proof search
is \emph{complete} with respect to
equational reasoning.
%
We define a notion of equational
proof
\eqpf{\Renv}{\Penv}{\expr}{\expr'}
and prove that
if there exists such a proof,
then
\pbe{\Renv}{\Penv}{\expr = \expr'}
is guaranteed to return \ttrue.
%
To prove this theorem, we introduce
the notion of \emph{bounded unfolding}
which corresponds to unfolding
definitions $n$ times.
%
We show that unfolding preserves
congruences, and hence, that an equational
proof exists \textit{iff} the goal can be proved with
\emph{some} bounded unfolding.
%
Thus, completeness follows by showing that
the proof search procedure computes the limit
(\ie fixpoint) of the bounded unfolding.
%
In \S~\ref{sec:pbe:terminates}
we show that the fixpoint is computable:
there exists an unfolding depth at which \pbesym
reaches a fixpoint and hence terminates.

\mypara{Bounded Unfolding}
%
For every $\Renv, \Penv$, and $0 \leq n$,
the \emph{bounded unfolding of depth} $n$ is
defined by:
%
$$\begin{array}{lcl}
  \binst{\Renv}{\Penv}{0}     & \doteq & \Penv \\
  \binst{\Renv}{\Penv}{n+1}   & \doteq & \Penv_n \cup \uinst{\Renv}{\Penv_n}
  \quad \mbox{where}\ \Penv_n = \binst{\Renv}{\Penv}{n}
\end{array}$$
%
That is, the unfolding at depth $n$ essentially
performs $\uinstsym$ upto $n$ times.
%
The bounded-unfoldings yield a monotonically
non-decreasing sequence of formulas such that
if two consecutive bounded unfoldings coincide,
then all subsequent unfoldings are the same.
%

\begin{lemma}[\textbf{Monotonicity}] \label{lem:bounded:mono}
$\forall 0 \leq n.\ \binst{\Renv}{\Penv}{n} \subseteq \binst{\Renv}{\Penv}{n + 1}$.
\end{lemma}

\begin{lemma}[\textbf{Fixpoint}] \label{lem:bounded:fix}
Let $\Penv_i \doteq \binst{\Renv}{\Penv}{i}$.
If $\Penv_n = \Penv_{n+1}$,
then $\forall n < m.\ \Penv_m = \Penv_n$.
\end{lemma}

\mypara{Uncovering}
%
Next, we prove that every function application
term that is \emph{uncovered} by unfolding to
depth $n$ is congruent to a term in the $n$-depth
unfolding.


\begin{lemma}[\textbf{Uncovering}] \label{lem:uncovering}
Let $\Penv_n \equiv \binst{\Renv}{\extendpenv{\Penv}{\vv = \expr}}{n}$.
If   $\smtIsValid{\Penv_n}{\vv = \expr'}$,
then for every $\issubterm{\fapp{\rfun}{t'}}{\expr'}$
       there exists $\issubterm{\fapp{\rfun}{t}}{\Penv_n}$
         such that $\smtIsValid{\Penv_n}{t_i = t'_i}$.
\end{lemma}

We prove the above lemma by induction on $n$
where the inductive step uses the following
property of congruence closure, which itself
is proved by induction on the structure of
$\expr'$:

\begin{lemma}[\textbf{Congruence}] \label{lem:congruence}
If   $\smtIsValid{\extendpenv{\Penv}{\vv = \expr}}{\vv = \expr'}$
     and
     $\vv \not \in \Penv, \expr, \expr'$,
then for every $\issubterm{\fapp{\rfun}{t'}}{\expr'}$
       there exists $\issubterm{\fapp{\rfun}{t}}{\Penv, \expr}$
         such that $\smtIsValid{\Penv}{t_i = t'_i}$.
\end{lemma}


\mypara{Unfolding Preserves Equational Links}
%
We use the uncovering Lemma~\ref{lem:uncovering}
and congruence to show that \emph{every instantiation}
that is valid after $n$ steps is subsumed by the $n+1$
depth unfolding.
%
That is, we show that every possible \emph{link} in
any equational chain can be proved equal to
the source expression via bounded unfolding.

\begin{lemma}[\textbf{Link}] \label{lem:link}
If   $\smtIsValid{\binst{\Renv}{\extendpenv{\Penv}{\vv = \expr}}{n}}{\vv = \expr'}$,
then $\smtIsValid{\binst{\Renv}{\extendpenv{\Penv}{\vv = \expr}}{n+1}}{\uinst{\Renv}{\extendpenv{\Penv}{\vv = \expr'}}}$.
\end{lemma}


\begin{figure}[t]
$$
\inference{
}{
	\eqpf{\Renv}{\Penv}{\expr}{\expr}
}[\rulename{Eq-Refl}]
$$

$$\inference
  { \eqpf{\Renv}{\Penv}{\expr}{\expr''}
    \quad
    \Penv' = \uinst{\Renv}{\extendpenv{\Penv}{\vv = \expr''}}
    \quad
    \smtIsValid{\Penv'}{\vv = \expr'}
  }
  {\eqpf{\Renv}{\Penv}{\expr}{\expr'}}
 [\rulename{Eq-Trans}]
$$

$$\inference
  { \eqpf{\Renv}{\Penv}{\expr_1}{\expr_1'}
    \quad
		\eqpf{\Renv}{\Penv}{\expr_2}{\expr_2'}
		\quad
    \smtIsValid{\Penv}{\expr_1' \binop \expr_2'}
  }
  {\symeqpf{\Renv}{\Penv}{\expr_1}{\binop}{\expr_2}}
  [\rulename{Eq-Proof}]
$$
\caption{\textbf{Equational Proofs:} rules for equational reasoning.}
\label{fig:equational-proof}
\figrule
\end{figure}

\mypara{Equational Proof}
%
Figure~\ref{fig:equational-proof} formalizes
our rules for equational reasoning.
%
Intuitively, there is an \emph{equational proof}
that $\expr_1 \binop \expr_2$ under $\Renv, \Penv$,
written by the judgment
%
$\symeqpf{\Renv}{\Penv}{\expr_1}{\binop}{\expr_2}$,
%
if by some sequence of repeated function unfoldings,
we can prove that $\expr_1$ and $\expr_2$ are
respectively equal to $\expr_1'$ and $\expr_2'$
such that,
%
$\smtIsValid{\Penv}{\expr_1' \binop \expr_2'}$
%
holds.
%
Our notion of equational proofs
adapts the idea of type level
computation used in TT-based
proof assistants to the setting of
SMT-based reasoning, via the
directional unfolding judgment
\eqpf{\Renv}{\Penv}{\expr}{\expr'}.
%
In the SMT-realm, the explicit
notion of a normal or canonical
form is converted to the implicit
notion of the equivalence classes
of the SMT solver's congruence
closure procedure (post-unfolding).

\mypara{Completeness of Bounded Unfolding}
%
We use the fact that unfolding preserves
equational links to show that bounded unfolding
is \emph{complete} for equational proofs.
%
That is, we prove by induction on the structure
of the equational proof that whenever there is
an \emph{equational proof} of $\expr = \expr'$,
there exists some bounded unfolding that
suffices to prove the equality.

\begin{lemma}\label{lem:bounded-completeness}
  If   $\eqpf{\Renv}{\Penv}{\expr}{\expr'}$,
  then $\exists 0 \leq n.\ \smtIsValid{\binst{\Renv}{\extendpenv{\Penv}{\vv = \expr}}{n}}{\vv = \expr'}$.
\end{lemma}



\mypara{\pbesym is a Fixpoint of Bounded Unfolding}
%
We show that the proof search procedure \pbesym computes
the least-fixpoint of the bounded unfolding and hence,
returns \ttrue \textit{iff} there exists \emph{some} unfolding
depth $n$ at which the goal can be proved.

\begin{lemma}[\textbf{Fixpoint}]\label{lem:fixpoint}
$\pbe{\Renv}{\Penv}{\expr = \expr'}$ iff
$\exists n.\ \smtIsValid{\binst{\Renv}{\extendpenv{\Penv}{\vv = \expr}}{n}}{\vv = \expr'}$.
\end{lemma}

The proof follows by observing that
$\pbe{\Renv}{\Penv}{\expr = \expr'}$
computes the \emph{least-fixpoint}
of the sequence
%
$\Penv_i \doteq \binst{\Renv}{\Penv}{i} \label{eq:fix:1}$.
%
Specifically, we prove by induction on $i$
that at each invocation of $\pbeloop{i}{\Penv_i}$
in Figure~\ref{fig:pbe}, $\Penv_i$ is equal to
$\binst{\Renv}{\extendpenv{\Penv}{\vv = \expr}}{i}$,
which then yields the result.

\mypara{Completeness of \pbesym}
%
Finally, we combine Lemmas~\ref{lem:fixpoint}
and~\ref{lem:link} to show that
\pbesym is complete, \ie if there is an equational proof
that $\expr \binop \expr'$ under $\Renv, \Penv$,
then \pbe{\Renv}{\Penv}{\expr \binop \expr'} returns $\ttrue$.

\begin{theorem}[\textbf{Completeness}] \label{thm:completeness}
  If \symeqpf{\Renv}{\Penv}{\expr}{\binop}{\expr'}
  then \pbe{\Renv}{\Penv}{\expr \binop \expr'} = \ttrue.
\end{theorem}

\subsection{\pbesym Terminates} \label{sec:pbe:terminates}


So far, we have shown that our proof search
procedure $\pbesym$ is both sound and complete.
%
Both of these are easy to achieve
simply by \emph{enumerating} all
possible instances and repeatedly
querying the SMT solver.
%
Such a monkeys-with-typewriters
approach is rather impractical: it
may never terminate.
%
Fortunately, we show that in
addition to being sound and complete
with respect to equational proofs,
if the hypotheses are transparent,
then our proof search procedure always
terminates.
%
Next, we describe transparency and
explain intuitively why \pbesym
terminates.
%
We then develop the formalism needed
to prove the termination
Theorem~\ref{thm:termination}.

\mypara{Transparency}
%
An environment $\env$ is \emph{inconsistent}
if $\smtIsValid{\tosmt{\env}}{\tfalse}$.
%
An environment \env is \emph{inhabited}
if there exists some $\store \in \embed{\env}$.
%
We say $\env$ is \emph{transparent}
if it is either inhabited \emph{or}
inconsistent.
%
As an example of a \emph{non-transparent}
$\Penv_0$ consider the predicate
$\kw{lenA xs} = 1 + \kw{lenB xs}$,
where \kw{lenA} and \kw{lenB} are both
identical definitions of the list length
function.
%
Clearly there is no $\store$ that causes
the above predicate to evaluate to $\ttrue$.
%
At the same time, the SMT solver cannot
(using the decidable, quantifier-free theories)
prove a contradiction as that requires
induction over \kw{xs}.
%
Thus, non-transparent environments are
somewhat pathological and, in practice,
we only invoke $\pbesym$ on transparent
environments.
%
Either the environment is inconsistent,
\eg when doing a proof-by-contradiction,
or \eg when doing a proof-by-case-analysis
we can easily find suitable concrete
values via random~\cite{Claessen00}
or SMT-guided generation~\cite{Seidel15}.

\mypara{Challenge: Connect Concrete and Logical Semantics}
%
As suggested by its name, the
\pbesym algorithm aims to lift the
notion of evaluation or computations
into the level of the refinement logic.
%
Thus, to prove termination, we must
connect the two different notions of
evaluation,
the \emph{concrete} (operational) semantics
and
the \emph{logical} semantics being used by \pbesym.
%
This connection is trickier than appears at first glance.
%
In the concrete realm totality ensures that every
reflected function $\rfun$ will terminate when
run on any \emph{individual} value $\val$.
%
However, in the logical realm, we are working with
\emph{infinite} sets of values, compactly represented
via logical constraints.
%
In other words, the logical realm can be viewed
(informally) as an \emph{abstract interpretation}
of the concrete semantics. We must carefully argue
that despite the \emph{approximation} introduced
by the logical abstraction, the abstract
interpretation will also terminate.

\mypara{Solution: Universal Abstract Interpretation}
%
We make this termination argument in three steps.
%
First, we formalize how \pbesym performs
computation at the logical level via
\emph{logical steps} and \emph{logical traces}.
We show (Lemma~\ref{lem:step-reductions})
that the logical steps form a so-called
\emph{universal} (or \emph{must})
abstraction of the concrete
semantics~\cite{CousotCousot77,CGL92}.
%
Second, we show that if \pbesym diverges,
it is because it creates a strictly
increasing infinite chain,
$\binst{\Renv}{\Penv}{0} \subset \binst{\Renv}{\Penv}{1} \ldots$
which corresponds to an \emph{infinite logical trace}.
%
Third, as the logical computation
is a universal abstraction we use
inhabitation to connect the two
realms, \ie to show that an
infinite logical trace
corresponds to an infinite
concrete trace.
%
The impossibility of the latter
must imply the impossibility of
the former, \ie \pbesym terminates.
Next, we formalize the above
to obtain Theorem~\ref{thm:termination}.




\mypara{Totality}
%
A function is \emph{total} when its evaluation
reduces to exactly one value.
%
The totality of \RRenv can and is checked
by refinement types (\S~\ref{sec:formalism}).
%
Hence, for brevity, in the sequel
we \emph{implicitly assume}
that $\RRenv$ is total under $\env$.
%
\begin{definition}[\textbf{Total}]\label{asm:total}
Let $\body \equiv \rdef{x}{\tosmt{\pred}}{\tosmt{e}}$.
$\body$ is \emph{total} under \env and \RRenv, if
forall $\store\in \embed{\env}$:
%
\begin{enumerate}
\item if $\rsatisfies{\RRenv}{\store}{\pred_i}$, then
$ \exists \val.\ \evalsto{{\apply{\store}{\refapply{\RRenv}{e_i}}}}{\val}$,
\item if $\rsatisfies{\RRenv}{\store}{\pred_i}$ and $\rsatisfies{\Renv}{\store}{\pred_j}$, then $i=j$, and
\item there exists an $i$ so that  $\rsatisfies{\RRenv}{\store}{\pred_i}$.
\end{enumerate}

$\RRenv$ is \emph{total} under \env,
if every $\body \in \tosmt{\RRenv}$
is total under \env and \RRenv.
%
\end{definition}

\mypara{Subterm Evaluation}
%
As the reflected functions are total,
the Church-Rosser theorem implies that
evaluation order is not important.
%
To prove termination, we require
an evaluation strategy, \eg CBV,
in which if a reflected
function's guard is satisfied,
then the evaluation of the
corresponding function body
requires evaluating
\emph{every subterm}
inside the body.
%
As $\defIf\cdot$ hoists
\kw{if}-expressions out
of the body and into the
top-level guards, the below
fact follows from the
properties of CBV:

\begin{lemma}\label{lem:subterm-eval}
Let $\body \equiv \rdef{x}{\tosmt{\pred}}{\tosmt{e}}$ and
    $\rfun \in \RRenv$.
%
For every \env, \RRenv, and $\store \in \embed{\env}$,
if $\rsatisfies{\RRenv}{\store}{\pred_i}$ and
   $\issubterm{\fapp{\rfun}{\tosmt{e'}}}{\tosmt{e_i}}$,
then
   $\evalsto{{\apply{\store}{\refapply{\RRenv}{e_i}}}}{\ctxapp{\ctx}{\rawapp{\rfun}{\apply{\store}{\refapply{\RRenv}{\params{e'}}}}}}$.
\end{lemma}

\mypara{Logical Step}
%
A pair
$\symstep{\fapp{\rfun}{\expr}}{\fapp{\rfun'}{\expr'}}$
is a
$\Renv,\Penv$-\emph{logical step} (abbrev. step),
if
%
\begin{itemize}
  \item $\rlookup{\Renv}{\rfun} \equiv \rdef{x}{\pred}{\body}$,
  \item $\smtIsValid{\Penv \wedge Q}{\pred_i}$ for some $(\Renv, \Penv)$-instance $Q$, and
  \item $\issubterm{\fapp{\rfun'}{\expr'}}{\SUBSTS{\body_i}{x}{\expr}}$.
\end{itemize}

\mypara{Steps and Reductions}
%
Next, using Lemmas~\ref{lem:subterm-eval},~\ref{lem:pbe:semantics},
and the definition of logical steps, we show that every
logical step corresponds to a \emph{sequence} of steps
in the concrete semantics:

\begin{lemma}[\textbf{Step-Reductions}]\label{lem:step-reductions}
If   $\symstep{\fapp{\rfun}{\tosmt{e}}}{\fapp{\rfun'}{\tosmt{e'}}}$
     is a logical step under $\tosmt{\RRenv}, \tosmt{\env}$
		 and
		 $\store\in \embed{\env}$,
then $\evalsto{{\fapp{\rfun}{\apply{\store}{\refapply{\RRenv}{e}}}}}
              {\ctxapp{\ctx}{\rawapp{f}{\apply{\store}{\refapply{\RRenv}{\params{e'}}}}}}$
for some context $\ctx$.
\end{lemma}

\mypara{Logical Trace}
%
A sequence
%
$\fapp{\rfun_0}{\expr_0},\fapp{\rfun_1}{\expr_1},\fapp{\rfun_2}{\expr_2},\ldots$
%
is a $\Renv,\Penv$-\emph{logical trace} (abbrev. trace),
if % $\params{\expr_0} \equiv \params{x_0}$ for some variables $\params{x_0}$, and
$\symstep{\fapp{\rfun_i}{\expr_i}}{\fapp{\rfun_{i+1}}{\expr_{i+1}}}$
is a $\Renv,\Penv$-step, for each $i$.
%
Our termination proof hinges upon the following
key result: inhabited environments only have
\emph{finite} logical traces.
%
We prove this result by contradiction.
%
Specifically, we show by Lemma~\ref{lem:step-reductions}
that an infinite $(\tosmt{\RRenv},\tosmt{\env})$-trace
combined with the fact that \env is inhabited
yields \textit{at least one}
\emph{infinite concrete trace},
which contradicts the totality assumption.
%
Hence, all the $(\tosmt{\RRenv},\tosmt{\env})$
logical traces must be finite.

\begin{theorem}[\textbf{Finite-Trace}]\label{thm:finite-trace}
If % $\RRenv$ is total under \env
   \ \env is inhabited,
then every $(\tosmt{\RRenv},\tosmt{\env})$-trace is finite.
\end{theorem}

\mypara{Ascending Chains and Traces}
%
If unfolding $\Renv, \Penv$ yields
an infinite chain $\Penv_0 \subset \ldots \subset \Penv_n \ldots$,
then $\Renv, \Penv$ has an infinite logical trace.
%
We construct the trace
by selecting at level $i$
(\ie in $\Penv_i$)
an application term
$\fapp{\rfun_i}{\expr_i}$
that was created by
unfolding an application
term at level $i-1$
(\ie in $\Penv_{i-1}$).

\begin{lemma} [\textbf{Ascending Chains}] \label{lem:chain}
Let $\Penv_i \doteq \binst{\Renv}{\Penv}{i}$.
If there exists an (infinite) ascending chain
$\Penv_0 \subset \ldots \subset \Penv_n \ldots$,
then there exists an (infinite) logical trace
$\fapp{\rfun_0}{\expr_0},\ldots,\fapp{\rfun_n}{\expr_n},\ldots$.
\end{lemma}

\mypara{Logical Evaluation Terminates}
%
Finally, we prove that the proof search
procedure $\pbesym$ terminates.
%
If $\pbesym$ loops forever,
there must be an infinite
strictly ascending chain of
unfoldings $\Penv_i$, and
hence, by Lemma~\ref{lem:chain},
an infinite logical trace, which,
by Theorem~\ref{thm:finite-trace},
is impossible.
%
\begin{theorem}[\textbf{Termination}] \label{thm:termination}
  If $\env$ is transparent,
  then $\pbe{\tosmt{\RRenv}}{\tosmt{\env}}{\pred}$ terminates.
\end{theorem}


% In practice, to check if \env is inhabited
% we need only find suitable concrete values
% via random~\cite{Claessen00} or SMT-guided
% generation~\cite{Seidel15}.
